\documentclass[10pt,a4paper]{amsart}
\usepackage[margin=19mm]{geometry}
\usepackage{amsmath,amssymb,mathtools}
\usepackage[scr=rsfs,cal=euler]{mathalfa}
\usepackage{xfrac}
\usepackage[unicode,hidelinks,bookmarks=false]{hyperref}
\newcommand{\ie}{i.e.\ }
\newcommand{\cf}{cf.\ }
\newcommand{\resp}{respectively\ }
\newcommand{\Subsection}{\subsection}
\providecommand{\nobreakdash}{\nobreak\hbox{-}}
\setlength{\emergencystretch}{2em}
\multlinegap0pt
\usepackage{bm}
\usepackage{bbm}
\usepackage{dsfont}
\usepackage{xspace}
\usepackage{cancel}
\usepackage{mathtools}
\usepackage{amsthm}
\makeatletter
\@ifundefined{lemma}{\newtheorem{lemma}{Lemma}}{}
\makeatother
\newcommand{\E}{\mathbb{E}}
\renewcommand{\d}{\mathrm{d}}
\title[Accuracy criterion for mean field approximations of Markov p]{Accuracy criterion for mean field approximations of Markov processes on hypergraphs (Lemma 2)}
\author{Illes Horvath, Daniel Keliger}
\date{Source: arXiv:2201.02041v4 (CC BY 4.0)}
\begin{document}
\maketitle

\noindent\emph{Source and licence.} Accuracy criterion for mean field approximations of Markov processes on hypergraphs. Version arXiv:2201.02041v4 (CC BY 4.0), distributed under the Creative Commons Attribution 4.0 International licence, as recorded in the arXiv licence metadata for this exact version and shown on the abstract page https://arxiv.org/abs/2201.02041v4. Full rights notice: licenses/arxiv-2201.02041-CC-BY-4.0.txt.\par\smallskip

\noindent\textit{Editorial scope.} Counted unit: the Lemma printed as Lemma 2 of the article (source0.tex lines 1291--1297, source label \texttt{c:1}), delivered together with the article's own complete original proof of it (source0.tex lines 1299--1336, one \texttt{proof} environment of 38 lines). No mathematics is written, repaired or re-derived: statement and proof are the article's own text line for line. Required context retained uncounted: none. Every definition, display and label of those lines is retained unchanged. Omitted from this package and not redistributed: 1--1290, 1298--1298, 1337--1786. Nothing inside a retained excerpt is omitted or altered: every retained body is delivered exactly as the article's lines are written, so no omission receipt of any kind accompanies this package. Citations: the retained excerpts carry no citation command at all, so no bibliography entry of the article is transcribed and no reference is redistributed. Reference adaptation: none was needed and none was made; every reference inside the delivered unit resolves inside it, so no unresolved reference or citation is produced. Printed slips kept literal: no printed word, symbol, hypothesis or number of the counted statement or of its proof was altered. Layout and class: the article's own class is \texttt{article} with packages including bbm, amsthm, inputenc, babel, color, amsmath, amssymb, yhmath, amsfonts, mathrsfs, enumitem, mathtools, dsfont, tikz; the wrapper is the lane's pinned \texttt{amsart} editorial wrapper at 10pt on A4, which restates the theorem-like environments the retained text needs and adds the article's own macro definitions that the retained text needs (\texttt{\textbackslash E}, \texttt{\textbackslash d}), with extra wrapper packages (bm, bbm, dsfont, xspace, cancel, mathtools). The delivered numbering is the unit's own. Assets: no retained span contains a figure environment, a graphics-inclusion command, a PDF-page inclusion, a table or a TikZ picture; no image file is copied into this package, the package declares no asset dependency and no image file is redistributed with it. Licence: version arXiv:2201.02041v4 of this article is distributed under the Creative Commons Attribution 4.0 International licence, as recorded in the arXiv licence metadata for this exact version and shown on the abstract page https://arxiv.org/abs/2201.02041v4. A full rights notice is delivered at licenses/arxiv-2201.02041-CC-BY-4.0.txt. Characters: every retained excerpt is pure ASCII, so the delivered engine, XeLaTeX with its default Latin Modern text font, reports no missing character.
	\begin{lemma}
		\label{c:1}
		Assume $M=1$. Let $\left( \omega_j \right)_{j\in [N]}$ be arbitrary non-negative weights. There exists a constant $\bar{C}=\bar{C}(q_{max})$ such that
		\begin{align*}
			\E \left( \sup_{0 \leq  \tau \leq t} \left | \sum_{j \in [N]} \omega_j \left(\hat{\xi}_{j,s}(\tau)-z_{j,s}(\tau) \right) \right|\right) \leq \bar{C}(1+t)\sqrt{\sum_{j \in [N]} \omega_j^2}.
		\end{align*}
	\end{lemma}

	\begin{proof} (Lemma \ref{c:1})
		\begin{align*}
			&\sum_{j \in [N]} \omega_j \left(\hat{\xi}_{j,s}(t)-z_{j,s}(t) \right)=\sum_{j \in [N]} \omega_j \left(\hat{\xi}_{j,s}(0)-z_{j,s}(0) \right) + \\
			&\sum_{s' \neq s}\sum_{j \in [N]}\omega_j \left[ \mathcal{N}_{j,ss'}\left(\mathcal{K}_{j,ss'}(t) \right) -\int_{0}^{t}q_{ss'} \left( \zeta_{j}(\tau)\right) \hat{\xi}_{j,s'}(\tau) \d \tau \right] \\
			&-\sum_{s' \neq s} \sum_{j \in [N]}\omega_j \left[ \mathcal{N}_{j,s's}\left(\mathcal{K}_{j,s's}(t) \right) -\int_{0}^{t}q_{s's} \left( \zeta_{i}(\tau)\right) \hat{\xi}_{j,s}(\tau) \d \tau \right]+ \\
			&\sum_{s'} \int_{0}^{t} \sum_{j \in [N]}\omega_j q_{ss'}(\zeta_j(\tau)) \left(\hat{\xi}_{j,s'}(\tau)-z_{j,s'}(\tau) \right) \d \tau
		\end{align*}
		
		Using the independence of $\hat{\xi}_{j,s}(t)$:
		\begin{align*}
			&\E \left[ \left| \sum_{j \in [N]} \omega_j \left(\hat{\xi}_{j,s}(0)-z_{j,s}(0) \right)\right| \right] \leq \mathbb{D} \left(\sum_{j \in [N]} \omega_j \hat{\xi}_{j,s}(0) \right)= \\
			&\sqrt{\sum_{j \in [N]}\omega_j^2\mathbb{D}^2 \left(\hat{\xi}_{j,s}(0) \right)} \leq \sqrt{\sum_{j \in [N]}\omega_j^2}.
		\end{align*}		
		
		For the remaining terms, we can use the same bound for all $s,s'$  pairs. Notice
		\begin{align*}
			\left|\mathcal{K}_{j,ss'}(t) \right|=\int_{0}^{t}q_{ss'} \left( \zeta_{j}(\tau)\right) \hat{\xi}_{j,s'}(\tau) \d \tau \leq q_{\textrm{max}}t
		\end{align*}
		and the fact that $\mathcal{N}_{j,ss'}\left(\mathcal{K}_{j,ss'}(t) \right) -\int_{0}^{t}q_{ss'} \left( \zeta_{j}(\tau)\right) \hat{\xi}_{j,s'}(\tau) \d \tau $ is a martingale. Using Doob's inequality,
		\begin{align*}
			&\E \left( \sup_{0 \leq \tau \leq t} \left| \sum_{j \in [N]}\omega_j \left[ \mathcal{N}_{j,ss'}\left(\mathcal{K}_{j,ss'}(\tau) \right) -\int_{0}^{\tau}q_{ss'} \left( \zeta_{j}(\tau')\right) \hat{\xi}_{j,s'}(\tau') \d \tau' \right] \right| \right) \leq \\
			&\left(\E\left( \sup_{0 \leq \tau \leq t} \left| \sum_{j \in [N]}\omega_j \left[ \mathcal{N}_{j,ss'}\left(\mathcal{K}_{j,ss'}(\tau) \right) -\int_{0}^{\tau}q_{ss'} \left( \zeta_{j}(\tau')\right) \hat{\xi}_{j,s'}(\tau') \d \tau' \right] \right|\right)^2   \right)^{1/2}\leq\\
			&2\left(\E\left( \left| \sum_{j \in [N]}\omega_j \left[ \mathcal{N}_{j,ss'}\left(\mathcal{K}_{j,ss'}(t) \right) -\int_{0}^{t}q_{ss'} \left( \zeta_{j}(\tau)\right) \hat{\xi}_{j,s'}(\tau) \d \tau \right] \right|\right)^2 \right)^{1/2}=\\
			& 2\left(\E \left\langle\sum_{j \in [N]}\omega_j \left[ \mathcal{N}_{j,ss'}\left(\mathcal{K}_{j,ss'}(\cdot) \right) -\int_{0}^{\cdot}q_{ss'} \left( \zeta_{j}(\tau)\right) \hat{\xi}_{j,s'}(\tau) \d \tau \right] \right\rangle_{\!\! t}\right)^{1/2}=\\
			& 2\left(\E\left\langle\sum_{j \in [N]}\omega_j \mathcal{N}_{j,ss'}\left(\mathcal{K}_{j,ss'}(\cdot) \right) \right\rangle_{\!\! t}\right)^{1/2}=
			2\left(\sum_{j \in [N]}\omega_j^2 \E\left\langle\mathcal{N}_{j,ss'}\left(\mathcal{K}_{j,ss'}(\cdot) \right) \right\rangle_{t}\right)^{1/2}\leq\\
			& 2\left(\sum_{j \in [N]}\omega_j^2 q_{\max}t\right)^{1/2}\leq
			2(1+ q_{\max}t)\sqrt{\sum_{j \in [N]}\omega_j^2},
		\end{align*}
		where $\langle\cdot\rangle_t$ denotes quadratic variation, and we used that the quadratic variation of the integral term is $0$, and the quadratic variation of a pure jump process is the total of the squared jumps (for the Poisson process, jump size is 1, and this reduces to the total number of jumps).
		
		As for the last term,
		\begin{align*}
			&\int_{0}^{t}\E \left[ \left| \sum_{j \in [N]} \omega_j q_{ss'}(\zeta_j(\tau)) \left(\hat{\xi}_{j,s'}(\tau)-z_{j,s'}(\tau) \right) \right| \right] \d \tau \leq \\
			&\int_{0}^{t} \mathbb{D} \left(\sum_{j \in [N]} \omega_j q_{ss'}(\zeta_j(\tau)) \hat{\xi}_{j,s'}(\tau)  \right)\d \tau =\\
			& \int_{0}^{t} \sqrt{\sum_{j \in [N]} \omega_{j}^2 q_{ss'}^2(\zeta_j(\tau)) \mathbb{D}^2 \left( \hat{\xi}_{j,s'}(\tau)\right)} \d \tau \leq q_{\textrm{max}}t \sqrt{\sum_{j \in [N]} \omega_j^2}.
		\end{align*}
	\end{proof}

\end{document}
